% TOP_PROBLEM: {"id": 30006769, "title": "Berman-Hartmanis Isomorphism Conjecture", "category_id": 15, "category": {"id": 15, "name": "computer_science", "display_name": "Computer Science", "description": "Computational complexity, algorithms, and theoretical CS.", "slug": "computer-science", "order_index": 15, "created_at": "2026-07-31T15:26:25.671Z"}, "status": "open"}
% FIRST_POSTED: 2026-09-25
% ENTRY_KIND: statement_only
% !TeX program = lualatex
% Definition and sources only; this file is not an LLM solution attempt.
\documentclass[11pt]{article}
\usepackage[margin=25mm]{geometry}
\usepackage{fontspec}
\setmainfont{Latin Modern Roman}
\usepackage{amsmath,amssymb,mathtools}
\usepackage{unicode-math}
\setmathfont{Latin Modern Math}
\usepackage{xurl,hyperref}
\hypersetup{colorlinks=true,urlcolor=blue}
\setcounter{secnumdepth}{0}
\setlength{\parindent}{0pt}
\setlength{\parskip}{5pt}
\setlength{\emergencystretch}{3em}
\begin{document}
\section*{\texorpdfstring{Berman-Hartmanis Isomorphism Conjecture}{Berman-Hartmanis Isomorphism Conjecture}}\label{30006769}
\subsection{Definitions and mathematical statement}
A language $A\subseteq\{0,1\}^*$ is NP-complete under polynomial-time many-one reductions if $A\in\mathrm{NP}$ and every NP language $L$ has a polynomial-time map $f$ with $x\in L\iff f(x)\in A$. A polynomial-time isomorphism between $A$ and $B$ is a bijection $h:\{0,1\}^*\to\{0,1\}^*$ satisfying
\[
 x\in A\iff h(x)\in B,
 \qquad h,h^{-1}\text{ polynomial-time computable}.
\]
The conjecture asserts such an $h$ for every pair of NP-complete languages. Reductions in both directions alone do not supply a bijection or an efficiently computable inverse.
\par\smallskip\textit{Formulation and concept references:} \href{https://www.cs.umd.edu/~gasarch/open/juris.pdf}{[S1]}.

\subsection{Short English statement}
Are all NP-complete languages the same up to a reversible, polynomial-time change of their binary encodings?
\subsection{Sources}
\begin{itemize}
\item[S1]Allender, Cai, Fortnow, Gasarch, Immerman, Kurtz, Royer, and Williams, Open Problems by or Inspired by Juris Hartmanis, ACM SIGACT News 53(4), 2022, section 5.3, printed p. 7.
\url{https://www.cs.umd.edu/~gasarch/open/juris.pdf}.
\textit{Formulation source.}
\end{itemize}
\end{document}
